\documentclass[12pt,a4paper]{article}
\usepackage[utf8]{inputenc}
\usepackage[T1]{fontenc}
\usepackage[french]{babel}
\usepackage{amsmath, amssymb}
\usepackage{graphicx}
\usepackage{geometry}
\geometry{margin=2cm}
\usepackage{enumitem}
\usepackage{float}
\usepackage{fancyhdr}
\usepackage{titlesec}
\pagestyle{fancy}

% Nettoyer les styles par défaut
\fancyhf{}

% En-têtes
\fancyhead[L]{Fiches d'exercices pour Maximilien (1)} % à gauche
\fancyhead[R]{Matteo Dubresson} % à droite

% Pieds de page
\fancyfoot[C]{\thepage} % numéro de page centré

%Titre et généralités
\title{Fiche d'exercices \\ pour Maximilien (1)}
\date{25/09/2025}
\author{}

%Espace entre les subsections
\titlespacing*{\subsection}{0pt}{5ex plus 0.5ex minus 0.2ex}{1.5ex plus 0.2ex}




\begin{document}

\maketitle

\subsection*{Exercice 1 : }

\[
(u_n) \begin{cases} 
u_0 = 1 \\ 
u_{n+1} = \dfrac{9}{6 - u_n}
\end{cases}
\]

Montrer par récurrence :  
\[0 < u_n < 3 \quad \forall \, n \in \mathbb{N}\]

\bigskip

\subsection*{Exercice 2 : }
Soit a et b deux réels. On considère la proposition suivante : si a + b est irrationnel, alors a ou b sont irrationnels.

\begin{enumerate}
    \item Quelle est la contraposée de cette proposition?
    \item Démontrer la proposition.
    \item Est-ce que la réciproque de cette proposition est toujours vraie?
\end{enumerate}

\bigskip

\subsection*{Exercice 3 : }
Montrer que 
\[
n! \geq 2^n \quad \forall \, n \geq 4
\]

\bigskip


\subsection*{Exercice 4 :}
 
\begin{enumerate}
    \item Soit $a \in \mathbb{R}$ tel que $a > 0$. Montrer par récurrence que pour tout $n \in \mathbb{N}$ :
    \[
    (1 + a)^n \ge 1 + n a
    \]
    \item En déduire la limite de la suite $(v_n)$ définie pour tout $n \in \mathbb{N}^*$ par $v_n = \left(\frac{4}{3}\right)^n$.
\end{enumerate}

\bigskip

\subsection*{Exercice 5 :}

On considère la suite $(x_n)$ définie par $x_0 = 1$, $x_1 = 3$ et pour tout $n \in \mathbb{N}$ :
\[
x_{n+2} = 4x_{n+1} - 3x_n
\]

\begin{enumerate}
    \item Calculer $x_2$ et $x_3$.
    \item Soient les propositions $A(n) : « \, x_n = 3^n \, »$ et $B(n) : « \, x_{n+1} = 3^{n+1} \, »$.
    \item Montrer par récurrence double (ou forte) que pour tout $n \in \mathbb{N}$, la proposition $(A(n) \text{ ET } B(n))$ est vraie.
\end{enumerate}


\bigskip


\subsection*{Exercice 6 : }

Soit $(u_n)$ la suite définie pour tout $n \in \mathbb{N}$ par :
\[
u_n = 3^{2n+1} + 2^{n+2}
\]

\begin{enumerate}
    \item Calculer $u_0$ et $u_1$.
    \item Montrer que pour tout $n \in \mathbb{N}$, $u_n$ est un multiple de $7$.
\end{enumerate}

\end{document}